\begin{examplebox}
	\textbf{\textcolor{CExample}{例 1（基础）}}

	\textbf{\textcolor{CExample}{题目：}}
	掷两枚均匀骰子，求点数之和不等于7的概率。

	\textbf{\textcolor{CExample}{解题步骤：}}
	\begin{description}[style=nextline, leftmargin=2.8em, labelsep=0.5em, itemsep=0.45em]
		\item[\textbf{第一步（找对立事件）：}] 设事件 $A$ 为“点数之和等于7”，则所求事件为 $\overline{A}$。
			\par\textbf{理由：} 利用对立关系将“不等于7”转化为“等于7”的逆。
		\item[\textbf{第二步（算反面概率）：}] 骰子点数组合共36种等可能结果。和为7的有6种，故
			\[
				\begin{aligned}
					P(A) &= \frac{6}{36} \\
					&= \frac{1}{6}
			\end{aligned}
			\]
			\par\textbf{理由：} 古典概型计数。
		\item[\textbf{第三步（代入公式）：}] 由对立事件概率公式得
			\[
				\begin{aligned}
					P(\overline{A}) &= 1 - P(A) \\
					&= 1 - \frac{1}{6} \\
					&= \frac{5}{6}
			\end{aligned}
			\]
			\par\textbf{理由：} 直接代入 $P(\overline{A})=1-P(A)$。
	\end{description}

	\textbf{\textcolor{CExample}{关键结论：}}
	\textbf{所求 $P(\overline{A})=\dfrac{5}{6}$}

	\medskip
	\textbf{\textcolor{CExample}{例 2（稍变形）}}

	\textbf{\textcolor{CExample}{题目：}}
	某人射击命中率为0.8，独立射击3次，求至少命中1次的概率。

	\textbf{\textcolor{CExample}{解题步骤：}}
	\begin{description}[style=nextline, leftmargin=2.8em, labelsep=0.5em, itemsep=0.45em]
		\item[\textbf{第一步（找对立事件）：}] 设事件 $A$ 为“至少命中1次”，则 $\overline{A}$ 为“3次都未命中”。
			\par\textbf{理由：} 至少命中1次的反面是全未命中。
		\item[\textbf{第二步（算反面概率）：}] 每次未命中概率 $1-0.8=0.2$，独立射击3次，故
			\[
				\begin{aligned}
					P(\overline{A}) &= 0.2^3 \\
					&= 0.008
			\end{aligned}
			\]
			\par\textbf{理由：} 独立性下概率乘法。
		\item[\textbf{第三步（代入公式）：}] 由 $P(A)=1-P(\overline{A})$ 得
			\[
				\begin{aligned}
					P(A) &= 1 - 0.008 \\
					&= 0.992
			\end{aligned}
			\]
			\par\textbf{理由：} 对立事件概率公式。
	\end{description}

	\textbf{\textcolor{CExample}{关键结论：}}
	\textbf{至少命中1次的概率为 $0.992$}

	\medskip
	\textbf{\textcolor{CExample}{例 3（综合应用）}}

	\textbf{\textcolor{CExample}{题目：}}
	袋中装有4个白球，3个黑球，从中不放回地取2球，求至少取到1个黑球的概率。

	\textbf{\textcolor{CExample}{解题步骤：}}
	\begin{description}[style=nextline, leftmargin=2.8em, labelsep=0.5em, itemsep=0.45em]
		\item[\textbf{第一步（找对立事件）：}] 设事件 $A$ 为“至少取到1个黑球”，则 $\overline{A}$ 为“取到的2个球都是白球”。
			\par\textbf{理由：} 对立事件定义。
		\item[\textbf{第二步（算反面概率）：}] 总的取法数 $\binom{7}{2}=21$，取2个白球的方法数 $\binom{4}{2}=6$，故
			\[
				\begin{aligned}
					P(\overline{A}) &= \frac{6}{21} \\
					&= \frac{2}{7}
			\end{aligned}
			\]
			\par\textbf{理由：} 古典概型计数。
		\item[\textbf{第三步（代入公式）：}] 由对立事件概率公式得
			\[
				\begin{aligned}
					P(A) &= 1 - P(\overline{A}) \\
					&= 1 - \frac{2}{7} \\
					&= \frac{5}{7}
			\end{aligned}
			\]
			\par\textbf{理由：} 直接代入 $P(A)=1-P(\overline{A})$。
	\end{description}

	\textbf{\textcolor{CExample}{关键结论：}}
	\textbf{至少取到1个黑球的概率为 $\dfrac{5}{7}$}
\end{examplebox}
